\section{What remains open}\label{sec:open}

The classification of \cref{thm:nf-intro} reduces the five-valent problem in general position to one
operator, $T=J_1\cdot(J_2\times J_3+J_4\times J_5)$, and what is not proved about it can
be stated precisely.

\begin{conjecture}[The equal-weight gap]\label{conj:B}
Every nonzero eigenvalue of $8T$ satisfies
\[
  q\ \ge\ 2\sqrt{18-2\sqrt{41}}=4.5579607\ldots,
\]
with equality only at $(\frac12;\frac12,1;1,1)$ up to the symmetries of \cref{prop:S4}.
\end{conjecture}

For $j_1=1$ the conjecture is proved, since \cref{thm:jone} gives
$q\ge4.558>2\sqrt{18-2\sqrt{41}}$, which also excludes equality on that slice. For
$j_1=2$, for all spins at most $19$ and for $j_1=\frac52$ with spins at most $35$ the
proofs give $q\ge4.5579607$, short of the conjectured constant by less than
$4\times10^{-8}$; a second certified pass at a threshold between $4.5579608$ and $5.56$
that excludes the minimiser would give the sharp form on the finite ranges. For
$j_1=\frac12$ and $\frac32$ only the weaker constants $2.70$ and $3.7$ are proved. With
\cref{thm:closing,thm:t0} the conjecture would complete the labelled dichotomy for
tangents in general position: such an embedding has a gap exactly when its $\sigma$ does
not contain $\pm4$.

\begin{conjecture}[The growth law]\label{conj:L}
There are constants $c_2,C>0$ such that every nonzero eigenvalue of $8T$ satisfies
\[
  q\ \ge\ \max\Bigl(4\sqrt{10-6\sqrt2}\;j_1,\ c_2\,\frac{M}{\ln M}\Bigr)-C,\qquad M=2j_{\max}.
\]
\end{conjecture}

Both terms are sharp in form: the slope in $j_1$ cannot be larger (\cref{thm:slope}) and
the growth in $M$ cannot be faster than $M/\ln M$ (\cref{thm:N1}). Numerically the lower
envelope in $M$ is about $1.5\,M/\ln M$. Along the two families of \cref{thm:N1},
$\qm\ln M/M$ is $3.151$ and $5.97$ at $M=2^{16}$, close to the heuristic values $\pi$ and
$2\pi$ of \cref{sec:large} but still rising, so its limits are not established. The linear part, with exactly this slope, is proved on $\Cl_6$ (\cref{thm:C6}); for
unrestricted tuples it is open. Behind it is a
statement about three spins alone: the smallest nonzero eigenvalue modulus $\nu$ of
$U_z$ on any weight space of any $j_2\otimes j_3\otimes j_4$ is at least
$\nu_0=\sqrt{(5-3\sqrt2)/2}$, attained only at $(\frac12,1,1)$ with weight $\pm\frac12$.
This holds for all triples with $2j\le12$ and is open in general. The second route to
large $j_1$, the pencil of \cref{thm:pencil} for triples of unbounded size, would need
the next-order structure of $\mathsf A$ rather than norm bounds.

\subsection{Equal weights at large \texorpdfstring{$j_1$}{j1}: a spin-\texorpdfstring{$j_1$}{j1} ladder}\label{subsec:ladder}

The rest of this section describes what we know about the open case,
$j_1\ge\frac52$, beyond the theorems of \cref{sec:tone}. None of it closes the case.

\paragraph{The reduced obligation.}
By \cref{lem:mu}, $q\ge2\sqrt\mu\,\lambda(W_0)$ with $2\sqrt\mu\ge4\sqrt{2s(s+1)/\pi}$, a
factor that grows linearly in $s=j_1$. So a much weaker statement about the compressed
Kronecker sum than a uniform gap would already suffice:
\begin{quote}
\textup{(RO)} There is $c>0$ such that every tuple with $j_1\ge\frac52$ that is not
covered by \cref{thm:tone} has a pairing with $\lambda(W_0)\ge c/(j_1+1)$.
\end{quote}
Then $q\ge4c\sqrt{2s/(\pi(s+1))}\ge2.697\,c$ on those tuples, and
\cref{conj:B} would hold with a weaker constant. The evidence for (RO) is strong. A
random search over $10\,000$ uncovered tuples with $5\le2j_1\le41$ and doubled leg spins
up to $200$, drawn from six regimes (nearly equal legs, two large legs, two small legs,
unbalanced pairs, large gaps, legs near $j_1$), found no counterexample: the smallest
$\lambda(W_0)$, in the best pairing, is $0.614$ for half-integer $j_1$ and $1.0033$ for
integer $j_1$, and the smallest $\lambda(W_0)(j_1+1)$ is $2.46$. Every half-integer minimiser
has four nearly equal legs, and along such families $\lambda(W_0)$ decreases slowly with the
legs, to $0.583$, $0.567$ and $0.549$ at $2j_1=41$, $81$ and $161$ with doubled legs up to
$2048$, consistent with a limit of $\frac12$. So the factor $1/(j_1+1)$ in (RO) is never
close to being needed; the data suggest $\lambda(W_0)\ge\frac12$.

\paragraph{The picture: the low modes are a spin-$j_1$ channel.}
One might expect the low singular vectors of $W_0$ to be edge states on the two walls
$d=\pm s$, where the cut acts. They are not. Along the extremal family
$(s;n-\frac12,n,n,n)$, the weight of the lowest mode on each diagonal matches, to within
$0.05$, the weights $|\langle m|m_x=\frac12\rangle|^2$ of the eigenvector of $J_x$ in spin
$s$ with eigenvalue $\frac12$, for $2s=5$, $9$ and $21$; along the antidiagonal $u$ the mode
is spread evenly over the whole strip. Keeping only one wall gives $\lambda(W_0)$ between
$0.90$ and $1.03$, and the interior strip $|d|\le s-1$ is only $6$--$30\%$ better than the
full one. The low modes are the $m_x=\pm\frac12$ channel of a spin-$s$ representation of
$\mathfrak{su}(2)$ acting across the diagonals, and they need both walls.

The $\mathfrak{su}(2)$ is visible in the operator itself. Between the diagonals $d$ and $d+1$ the
factor $\sqrt{F(d+\frac12)}=2\ell_s(d)$ of every link is twice the link
$\ell_s(d)=\frac12\sqrt{(s-d)(s+d+1)}$ of $J_x$ in spin $s$, and \cref{prop:grid} gives at
once:

\begin{proposition}{$T$ is a $J_x(s)$ ladder with a diagonal defect}{Tladder}
For every tuple and the pairing $(23|45)$, every link of $T$ between the diagonals $d$
and $d+1$ is $\ell_s(d)\,\nu$, with $\nu=\frac12h_1(b)\sqrt{G(u)}$ for an $X$-link and
$\nu=\frac12h_2(c)\sqrt{G(u)}$ for a $D$-link, $(b,c)$ its upper end. So
$T=\sum_d\ell_s(d)(N_d+N_d^T)$, with $N_d$ the block of $\nu$-links from diagonal $d$ to
$d+1$, and on every interior diagonal $N_d^TN_d-N_{d-1}N_{d-1}^T$ is diagonal, with entry
$\frac14[G(u+1)P(b+1,c+1)-G(u)P(b,c)]$ at $(b,c)$, where $P(b,c)=h_1(b)^2-h_2(c)^2$.
\end{proposition}

If that diagonal defect vanished, $T$ would be exactly $J_x(s)\otimes P\oplus0$ for some
$P\succeq0$, by the argument of \cref{thm:ladderA} below with the Gram matrices
$N_d^TN_d$ in place of its Casimir (\cref{app:ladder}), and its smallest nonzero
eigenvalue would be $\frac12$ or $1$ times the smallest nonzero eigenvalue of $P$. In
the frame of $W_0$ the same structure appears with an absolute constant: the plaquette
argument of \cref{prop:zigzag} gives, for $W_0$, $M_d^TM_d-M_{d-1}M_{d-1}^T=-\frac d2+\Delta_d$
on every interior diagonal, with $M_d$ the block of $W_0$ from diagonal $d$ to $d+1$ and
$\Delta_d$ diagonal; and $-\frac d2$ is exactly the relation
$\ell_s(d)^2-\ell_s(d-1)^2=-\frac d2$ of $J_x$. With no defect this structure forces a
gap, uniformly in the sizes of the blocks:

\begin{theorem}{The su(2) ladder without defect}{ladderA}
Let $V_{-s},\dots,V_s$ be finite-dimensional real inner-product spaces, $M_d:V_d\to V_{d+1}$
linear, and $W=\sum_d(M_d+M_d^T)$ on $\bigoplus_dV_d$. Suppose
$M_d^TM_d-M_{d-1}M_{d-1}^T=-\frac d2$ on $V_d$ for every $-s<d<s$, and
$M_{s-1}M_{s-1}^T\succeq\frac s2$ on $V_s$, $M_{-s}^TM_{-s}\succeq\frac s2$ on $V_{-s}$. Then
$W$ is an orthogonal sum of copies of $J_x(S)$ with $S\le s-1$, $S\equiv s\pmod1$, and of
the chains $J(\kappa)$ on $-s\le d\le s$ with links $\sqrt{\kappa-d(d+1)/4}$ and
$\kappa\ge s(s+1)/4$. Every singular value of the bipartite block of $J(\kappa)$ is
nondecreasing in $\kappa$, and $J(s(s+1)/4)=J_x(s)$. Consequently every eigenvalue of $W$
has modulus at least $\frac12$ if $s$ is a half-integer, and every nonzero eigenvalue has
modulus at least $1$ if $s$ is an integer. Both bounds are sharp.
\end{theorem}

The proof (\cref{app:ladder}) is algebra: $C_d=M_d^TM_d+\frac{d(d+1)}4$ is intertwined by
the blocks, $M_dC_d=C_{d+1}M_d$, so $W$ commutes with a Casimir and splits into its
eigenspaces, on which every block is a scalar multiple of an isometry. The constants
$\frac12$ and $1$ are the bulk limits seen in the data.

\paragraph{Why the theorem does not apply.}
For the real $W_0$ the defect is
\[
  \Delta=\tfrac12\bigl(\varepsilon_1(b)-\varepsilon_2(c)\bigr)+[u=s]\bigl(h_1(b)^2-h_2(c)^2\bigr),
  \qquad \varepsilon(L)=\frac{(4n^2-1)(4m^2-1)}{2(2L-1)(2L+1)(2L+3)},
\]
with $\varepsilon_1$, $\varepsilon_2$ formed with the parameters of the two pairs, and with
the convention $\varepsilon(\frac12)=(4n^2-1)/8$ at the bottom label of a pair with
$m=\frac12$; the last term is present only at the corner sites $u=s$. It vanishes only
for two pairs with $m=\frac12$ away from their bottom labels and from the corner, which
no real tuple achieves everywhere. And it is not small. On eight tuples near the
minimisers its largest entry is $112$ to $1\,110$ in the band $u\le2s+2$ at the corner,
against at most $53$ elsewhere, and its size relative to the Gram matrices is between
$0.74$ and $11$. On the tuples with nearly equal legs the lowest mode's Casimir has its
mean within about one sector of the bottom value $s(s+1)/4$, but a spread of several
sectors that grows with the legs. In the
frame of $T$ the defect of \cref{prop:Tladder} is at most $9.4\%$ of the Gram diagonal
entry by entry on the two tuples measured, but it is still of order one as an operator, because the small
singular values of each $N_d$ come from cancellation along $u$. So a perturbation of
\cref{thm:ladderA} cannot reach the real operator.

\paragraph{What fails, and what a proof needs.}
Two simpler routes are false. The gap is not monotone in $j_1$ with legs $2$--$5$ fixed:
$q=148.70$ at $(4;1,\frac{11}2,\frac{11}2,6)$ but $70.59$ at $(5;1,\frac{11}2,\frac{11}2,6)$,
and $912$ of $10\,252$ such steps without a leg $\frac12$ go down. Nor is it monotone along
$j_1-j_5$ fixed, the setting of \cref{thm:pencil}: $q=146.46$ at $(5;1,4,4,1)$ but
$63.30$ at $(\frac{11}2;1,4,4,\frac32)$, with $490$ of $19\,498$ steps going down. Reading the
strip along the antidiagonals instead of the diagonals gives a transfer-matrix version of
\cref{blem:evenchain}, but its bounds collapse on unbalanced pairs, for example to
$1.3\times10^{-26}$ where $q=639$.

A proof of \cref{conj:B} along the lines of this section needs a version of
\cref{thm:ladderA} that allows the diagonal defect $\Delta$, under a hypothesis that is
structural rather than a bound on its size, that the real $W_0$ satisfies in some
pairing, and that still yields $\lambda(W_0)\ge c/(j_1+1)$. The defect is explicit and
concentrated at the corner of the strip, so the hypothesis presumably says that the
corner acts as a trivial boundary condition, the two-dimensional counterpart of the
trivial dimerisation of the $j_1=\frac12$ chain. We have not found it.

\paragraph{Identities behind this section.}
These are proved, and each was checked against an independent computation; the
scripts and outputs are in \texttt{certificates/su2l}, \texttt{w0ob} and \texttt{t1all}.
\emph{Lemma O}: on the invariant space $T=\frac1{2i}[J_2\cdot J_3+J_4\cdot J_5,\Omega]$ with
$\Omega=(J_3+J_5)^2-(J_2+J_4)^2$ (residual $5\times10^{-15}$ by brute force on all $512$
tuples with doubled spins at most $4$). \emph{Lemma U}: along the antidiagonals the Gram
matrices of $T$ satisfy the analogue of \cref{prop:zigzag} with ratio $G(u+1)/G(u)$ (exact,
$97\,597$ entries). \emph{Lemma S}: for a Stark matrix, $h(L+1)^2-h(L)^2=
-\frac{2L+1}4+\frac12\varepsilon(L)$, and for $m=\frac12$ the Stark matrix is $J_x$ of spin
$n-\frac12$ on its odd subspace (exact, $20\,540$ labels). \emph{Lemma L}: the Gram identity
for $W_0$ with the defect $\Delta$ above (exact, $195\,361$ sites, $27\,993$ at the corner).
\emph{Lemma W}: for half-integer $j_1$, in every pairing exactly one wall has the corner's
colour, and the leak through the other wall obeys the same $J_x$ law (exact, $12\,900$
tuples). \emph{Lemma Q}: an explicit two-spin Casimir commutes with $K$ up to the
$\varepsilon$-terms (float, $16\,875$ pairs of pairs). \Cref{prop:Tladder} was checked
exactly on all $55\,000$ tuples with $2j_1\le11$ and doubled legs at most $10$, and
\cref{thm:ladderA} on $3\,312$ honest two-spin strips and $400$ random exact ladders.

\paragraph{Beyond valence five.}
At valence six we have one embedding class (\cref{thm:v6}), chosen because Brunnemann and
Rideout's data single it out; its $\sigma$ has only two nonzero entries, and its mechanism
is close to the five-valent one. Whether the classes that Brunnemann and Rideout found
increasing at valence six, including the cubic vertex, have a gap is open, as is
everything non-degenerate at valence seven and higher.

\paragraph{Other operators.}
The general five-valent gap of the Rovelli--Smolin operator is open, and nothing here
concerns the polyhedral volume of Bianchi, Don\`a and Speziale~\cite{BDS2011}, for which
Haggard's level-repulsion argument suggests a generic gap~\cite{Haggard2013}. The
semiclassical status of the small Ashtekar--Lewandowski eigenvalues is itself debated: the
operator does not reproduce the polyhedral volume at generic five-valent
vertices~\cite{BDS2011,FT2008}, and the gap closing found here may be read as a property
of its regularisation rather than of quantum geometry. It may still matter for
dynamics. Brunnemann and Thiemann noted that the spectrum near zero volume enters the
Hamiltonian constraint through the curvature operator, and that valleys of zero
volume could make the relevant differences of the spectrum unbounded~\cite[p.~46]{BT2006};
eigenvalues as small as those of \cref{thm:closing} and \cref{thm:v6} sharpen that concern
for the embeddings that carry them. The classification says which embeddings carry the
effect; whether they should be counted is a question about the dynamics that this paper
does not address.
